\section{TLS}

TLS (\emph{Transport Layer Security}) is a protocol used to secure data
exchange between two hosts. It is mainly known for HTTPS, which is HTTP
over TLS. TLS provides authentication of the communicating hosts,
confidentiality through encryption, and integrity protection for the
transmitted data.

\paragraph{TLS Handshake}

Before encrypted communication can start, the client and server perform a
handshake. In TLS 1.2, the client first opens a TCP connection and sends
a \texttt{ClientHello} containing a random value $R_C$ and the list of
supported cryptographic algorithms. The server replies with its own
random value $R_S$, the chosen algorithms, and its X.509 certificate.

The client then verifies the server certificate. This allows the client
to check that it is really communicating with the intended server and not
with a Man-in-the-Middle.

$R_C$ and $R_S$ are fresh public random values used in the TLS handshake
to derive unique session keys and prevent replay of old handshakes.
\paragraph{Certificates}

A certificate is a signed statement binding a server identity
(e.g. domain name) to a public key. It is signed by a Certificate
Authority (CA) using the CA's private key. When the client receives the
certificate, it verifies the CA's signature using the CA's public key,
which is trusted through the browser/OS root certificate store.

\paragraph{Chain of Trust}

A client can verify a certificate only if it trusts the CA that signed
it. In practice, certificates are often signed by intermediate CAs, which
are themselves signed by higher-level CAs. This forms a \emph{chain of
trust}.


\paragraph{Key Exchange}

After the certificate has been verified, client and server establish
shared secrets. In TLS 1.2, this can be done using Diffie-Hellman.

The server sends its DH parameters and signs them with its private key (no men in the middle !).
The client verifies the signature using the public key from the
certificate. Then both sides compute the same Premaster Secret:

\[
PS = g^{ab} \bmod p
\]

From the Premaster Secret and the random values $R_C$ and $R_S$, both
sides derive a Master Secret using a pseudo-random function. From this,
TLS derives symmetric encryption keys and MAC keys, usually one pair for
each direction.

\paragraph{TLS Records}

After the handshake, user data is sent in TLS records. A TLS record
contains a fragment of application data, usually up to 16 KB. In TLS 1.2,
the record is typically authenticated with a MAC and then encrypted with
the negotiated symmetric cipher.

Sequence numbers are included in the MAC computation to prevent replay
attacks.

% \paragraph{RSA Key Exchange}

% Older TLS versions could also use RSA for key exchange. In that case,
% the client generates a Premaster Secret and encrypts it with the
% server's public RSA key.

% This is no longer recommended because it does not provide forward
% secrecy: if the server's long-term private key is compromised later, old
% recorded sessions may be decrypted.

\paragraph{Public Key Infrastructure}

The certificates, Certificate Authorities, revocation lists, OCSP
servers, and related components form a Public Key Infrastructure (PKI).
The PKI is what allows a browser to decide whether a server certificate
should be trusted.

\paragraph{Limitations of TLS}

TLS protects the communication channel, but it does not solve every
security problem. It does not prevent DoS attacks, SQL injection, XSS, or
phishing. It also cannot protect users if they trust the wrong website or
if a trusted CA is compromised.

\paragraph{TLS 1.3}

TLS 1.3 improves TLS 1.2 by removing weak algorithms and problematic
modes. RSA is no longer allowed for establishing the Premaster Secret;
instead, TLS 1.3 uses DH or ECDH, which provide forward secrecy.

TLS 1.3 also signs the cipher negotiation, which helps prevent downgrade
attacks where a Man-in-the-Middle tries to force client and server to use
weaker cryptography.

\paragraph{AEAD Ciphers}

TLS 1.3 only allows AEAD ciphers which stands for \emph{Authenticated Encryption with
Associated Data}. It performs encryption and authentication together and
produces both a ciphertext and an authentication tag.

This avoids the older separate MAC-then-encrypt construction.

\paragraph{Latency Improvements}

TLS 1.3 reduces handshake latency. In 1-RTT mode, the client already
sends key exchange information in its first message, allowing encrypted
communication to start after one round trip.

TLS 1.3 also supports 0-RTT mode for resumed connections, where the
client reuses a previously established secret. However, 0-RTT data is
vulnerable to replay attacks, so it should not be used for sensitive
operations such as bank transfers.

\paragraph{Server Name Indication}

During the TLS handshake, the client sends the server name it wants to
connect to using Server Name Indication (SNI). This is necessary because
one IP address may host several domains with different certificates.

In TLS 1.2 and normal TLS 1.3, SNI is sent in clear text, so observers
can see which website is being visited. Encrypted Client Hello (ECH) is a
TLS 1.3 extension that aims to encrypt this information.
